\documentclass[12pt]{article}
\usepackage[utf8]{inputenc}

%%% I added a few packages. This next one allows you to set the margins.
\usepackage[top=.5in,bottom=1in,right=.75in,left=.75in]{geometry}
\usepackage{multicol}
%\usepackage{newspaper}
\setlength{\columnsep}{1cm}
\usepackage{blindtext}
\usepackage{varwidth}

%%% This one does boxes in a really nice way.
\usepackage{tcolorbox}
\usepackage{mdframed}
\usepackage{minibox}
\usepackage{graphicx}
\graphicspath{ {./images/} }
\usepackage{epstopdf}
  \epstopdfDeclareGraphicsRule{.pdf}{png}{.png}{convert #1 \OutputFile}
  \DeclareGraphicsExtensions{.png,.pdf}
\usepackage{hyperref}
\hypersetup{
    colorlinks=true,
    linkcolor=blue,
    filecolor=magenta,      
    urlcolor=cyan,
    pdftitle={MSCSMess},
    pdfpagemode=FullScreen,
    }
\urlstyle{same}
 \usepackage{fix-cm}
 \usepackage{soul}
 \usepackage{graphicx}
 \usepackage{multicol}


\newcommand{\ds}{\displaystyle}
\newcommand{\ZZ}{\mathbb{Z}}
\newcommand{\QQ}{\mathbb{Q}}
\newcommand{\RR}{\mathbb{R}}
\newcommand{\CC}{\mathbb{C}}
\newcommand{\FF}{\mathbb{F}}
\newcommand{\xx}{\mathbf{x}}
\newcommand{\uu}{\mathbf{u}}
\newcommand{\vv}{\mathbf{v}}

\newcommand{\bolder}{\textbf}
 
 
\title{\fontsize{60}{70}\selectfont {\bf {MSCS} 
\includegraphics[scale=.3]{images/OleLion.png} {Mess}}}
\author{Department of Mathematics, Statistics, and Computer Science \\ St.\@ Olaf College, Northfield, MN 55057 \vspace{-2ex}} 
\date{\today\, $|$ Volume 54, No.\@ 9 \line(1,0){500}}


\begin{document}

\maketitle \vspace{-0.5cm}

Happy Friday! We hope that your registration for J-term and spring classes is going well. Check out the events that are happening next week and have a great weekend!


 \begin{multicols}{2}

\begin{tcolorbox}[colback=white, colframe=black]
    \begin{center}
    \title{{\large \bolder{ NSM Action Team Event: Pathways }}}
    \end{center}
\end{tcolorbox}
 Are you finding your STEM courses challenging and wondering what to expect next year? Are you thinking about what courses to take or what internships and research to get involved in? Are you interested in sharing your own experiences with other students? Come join the Natural Science and Mathematics Student ACTION Team, next Thursday, Nov 13,  from 5-6pm for pizza in Regents 356. Hear how senior Biology, Chem, Physics, Psych, Computer Science, Math and Stats majors have 
 navigated their STEM journey! 

\begin{tcolorbox}[colback=white, colframe=black]
    \begin{center}
    \title{{\large \bolder{Women Mathematicians in Quant Finance}}}
    \end{center}
\end{tcolorbox}

If you’re an undergrad student who is curious about how your math skills can be leveraged in different career paths, this session is for you! Join the Carnegie Mellon University Master of Science in Computational Finance (MSCF) program for a thought-provoking webinar that introduces the world of quantitative finance—a field where math, data, and real-world impact intersect. Hear from current female students in the Carnegie Mellon University MSCF program, alongside a successful alumna moderator, as they share why they chose this interdisciplinary field and how it’s opening doors to exciting, high-impact careers. Whether you have heard of quant finance or are just starting to explore your options, this session will help you see how your math skills can lead to opportunities beyond academia. Join the Association of Women in Mathematics for "Shifting the Curve: Women Mathematicians in Quant Finance" on November 13th, 6:00-7:30pm, EST. Registration is required, see poster below for details.



\begin{tcolorbox}[colback=white, colframe=black]
    \begin{center}
    \title{{\large \bolder{Statistics and Data Science Club!
}}}
    \end{center}
\end{tcolorbox}

Hello everyone! The Statistics and Data Science Club is back, and they would love to have you join! Apply your data skills to real-world projects, explore mentorship opportunities, and grow your expertise beyond the classroom. Sign up for the email list \href{https://docs.google.com/forms/d/e/1FAIpQLSeV00ldRadWwtvfA8nmOIuL43iZ0MFHp6O5cmrPY5HAsxVJMQ/viewform}{here} to get updates about meetings and events. 


\begin{tcolorbox}[colback=white, colframe=black]
    \begin{center}
    \title{{\large \bolder{MSCS student profile in the St. Olaf Magazine
}}}
    \end{center}
\end{tcolorbox}

There is an excellent article about Jessica Schmidt--math/bio major, SDS/math bio concentrator, class of 2026--in the current issue of the \href{https://wp.stolaf.edu/magazine/}{St. Olaf Magazine}, titled "The Cipher Protocol." Jessica is a "master codebreaker" after having been the only individual to have solved all three codebreaking challenges in the annual Kryptos competition. Read the article (and the follow-up article on the MSCS department and Math program) and congratulate Jessica.



\begin{tcolorbox}[colback=white, colframe=black]
    \begin{center}
    \title{{\large \bolder{ CS Unplugged   }}}
    \end{center}
\end{tcolorbox}

This year the CS program is hosting CS Unplugged: a four-part series on internships, jobs, and life as a computer scientist. These sessions will be informal, student-centered, and are a part of a larger goal to get CS students building community within the program. Fill out the form \href{https://docs.google.com/forms/d/e/1FAIpQLScfAnexFhdd-ZAHY0RDX8MpB6A7D7cFtjFYITs-VHnXKzg8XQ/viewform}{at this link} to give your input on what you want to hear at these events!

The first two events will be held on Wednesday, Nov 12th from 3:30-4:30pm, and Wednesday, Nov 19th from 3:30-4:30pm, both in RNS 203. At the end of each event, prizes will be given out! See the poster below for more info.

\begin{tcolorbox}[colback=white, colframe=black]
    \begin{center}
    \title{{\large \bolder{Problem Solving Group}}}
    \end{center}
\end{tcolorbox}
The Math Problem Solving Group meets weekly to have fun, solve math problems, eat candy, and prepare for upcoming math contests. Meetings are 11:30am-12:30pm, in RNS 207 (the Center for Interdisciplinary Research room). Everyone is welcome, regardless of prior experience in mathematics or math contests.

The first contest will be the Jerry Heuer Team Contest on the morning of Saturday, November 15. It's a fun, accessible contest for teams of up to three students. For more information, or to be added to the problem solving group email list, contact Prof. Wright at wright5@stolaf.edu.


\begin{tcolorbox}[colback=white, colframe=black]
    \begin{center}
    \title{{\large \bolder{DAAD RISE Internships}}}
    \end{center}
\end{tcolorbox}

Are you interested in a summer STEM internship/fellowship program for undergraduates in Germany, requiring no German-language skills? Check out the DAAD Research Internships in Science and Engineering (RISE) website \href{https://www.daad.de/rise/en/}{at this link}. The application window is October 15th - November 30th.

\end{multicols} \vspace{0.5cm}

\vfill

\begin{center}
\includegraphics[width=0.9\textwidth]{images/NSM action team Nov 13.png}
\end{center}

\begin{center}
\includegraphics[width=0.9\textwidth]{images/CS Unplugged.png}
\end{center}

\begin{center}
\includegraphics[width=0.9\textwidth]{images/AWMQuantFinance.png}
\end{center}

\begin{center}
\includegraphics[width=0.9\textwidth]{images/Tea Time Poster.png}
\end{center}





\noindent{To submit an article, event, or anything else for publication in the Mess, email messeditor@stolaf.edu. Visit \href{https://wp.stolaf.edu/mscs/mscs-mess/}{http://wp.stolaf.edu/mscs/mscs-mess/} for a digital archive of previous MSCS Mess issues.}\\
	
\hfill Nima Dahir, Editor
	
\hfill Kyle Salois, Faculty Advisor


 

\end{document}